\clearpage
\section{Binding-free paths and the boundary of plugging closure}
\label{sec:copying-binding}
This section treats the plugging-closure prerequisite of Appendix~B.1,
item~3. It does not complete that item: the copying observations and the
translation/congruence theorem for lambda-contexts remain unfinished.
We retain Definition~11.2 literally, prove the positive-path correspondence,
and prove Lemma~E.1 with a sourced-hole hypothesis. That hypothesis holds
for every translated term hole. A finite counterexample shows why it cannot
be dropped for arbitrary hole interfaces allowed by the current graph model.
Neither program of Proposition~12.2 nor its Figure~16 translation changes.

\subsection{The literal path restriction}
\termdef{copying-binding-free}
A binding-free graph is focus-free and has no positive path, at any depth,
from a source of a contraction, instance, box, or hole edge to a source of
a hole. A store edge is not an instance edge. Paths remain at one depth;
the condition is checked separately in each box body.


The binding-label predicate classifies the four prohibited starting kinds;
the hole-label predicate recognises an actual primitive hole. Reflexive
edge ancestry follows zero or more incidence adjacencies. Strict ancestry
requires a first adjacency and hence at least one traversed edge before the
terminal hole. The shallow restriction additionally requires a source of
the starting edge. This is substantive for holes with an empty input list.
The recursive path predicate conjoins the shallow restriction with the
same predicate in every box body. Finally, binding-freedom adds global
focus-freedom. These predicates do not themselves enforce unique hole names;
that independent context condition must also be supplied where required.

\lean{Copying/BindingFree}{11}{19}{BindingFree.BindingLabel}


\lean{Copying/BindingFree}{20}{24}{BindingFree.HoleLabel}


\lean{Copying/BindingFree}{25}{28}{BindingFree.Ancestor}


\lean{Copying/BindingFree}{29}{33}{BindingFree.StrictAncestor}


\lean{Copying/BindingFree}{34}{38}{BindingFree.ShallowBindingFree}


\lean{Copying/BindingFree}{39}{44}{BindingFree.BindingPaths}


\lean{Copying/BindingFree}{45}{49}{BindingFree.BindingFree}


Ancestry is transitive by induction on its first derivation. An ancestry
between distinct edges must have a first step, since its reflexive case
would identify the endpoints. Any adjacency-preserving edge map transports
ancestry by induction. These facts concern incidence alone, without an
operational or termination hypothesis.

\lean{Copying/BindingFree}{50}{55}{BindingFree.trans}


\lean{Copying/BindingFree}{56}{61}{BindingFree.positive_of_ne}


\lean{Copying/BindingFree}{62}{72}{BindingFree.map}


A vertex step traverses a single edge from one of its sources to one of
its targets. Starting at any source of an ancestor edge, induction on
ancestry yields a finite vertex trace to a source of the terminal edge.
Conversely, a vertex trace beginning at a target of an edge yields strict
ancestry to any edge having the final vertex as a source. The induction
prepends each traversed edge.

Consequently, strict ancestry together with a starting-source witness is
\emph{equivalent} to a vertex trace of length $n+1$ between the two source
lists, for some $n$. In the reverse implication, source uniqueness identifies
the first traversed edge with the proposed starting edge. The positive
length is essential: reflexive vertex reachability would wrongly exclude
an isolated hole from binding-free contexts.

\lean{Copying/BindingPathCorrespondence}{7}{11}{BindingPathCorrespondence.BindingPathStep}


\lean{Copying/BindingPathCorrespondence}{12}{21}{BindingPathCorrespondence.ancestor_vertex_trace}


\lean{Copying/BindingPathCorrespondence}{22}{34}{BindingPathCorrespondence.target_trace_strict}


\lean{Copying/BindingPathCorrespondence}{35}{52}{BindingPathCorrespondence.strict_iff_vertex_path}


Every hole is a prohibited starting label. A self-loop at a hole would
supply both a source witness and a positive path back to that hole, so a
binding-free hole cannot self-loop. Globally hole-free graphs satisfy the
path restriction vacuously, first at the outer level and then by induction
inside boxes. Combining this with focus-freedom proves binding-freedom.
A general pullback lemma transports the shallow restriction along an edge
map preserving adjacency, starting-label classification, source existence,
and terminal-hole classification: map any proposed counterexample and
contradict the target graph's restriction.

\lean{Copying/BindingFree}{73}{79}{BindingFree.hole_binding}


\lean{Copying/BindingFree}{80}{85}{BindingFree.hole_nonSelfLoop}


\lean{Copying/BindingFree}{86}{94}{BindingFree.shallow_of_holeFree}


\lean{Copying/BindingFree}{95}{109}{BindingFree.paths_of_holeFree}


\lean{Copying/BindingFree}{110}{114}{BindingFree.of_holeFree}


\lean{Copying/BindingFree}{115}{128}{BindingFree.shallow_of_map}


\subsection{Plugging closure at sourced holes}
\termdef{copying-sourced-hole}
A sourced hole is an addressed hole occurrence whose input list has at
least one port. The condition concerns only the replaced hole, including
when that hole lies inside boxes. It does not exclude unrelated zero-input
holes from the surrounding observation context.


For an adjacency from a surviving host edge into an inserted edge,
source incidence of the inserted graph forces the shared vertex to be an
input port. Thus the old host edge was adjacent to the replaced hole.
Dually, an inserted-to-host adjacency comes from an output port and makes
the old hole adjacent to the surviving host edge.

Induction backwards along an ancestry ending at a hole now identifies its
origin. A surviving starting edge has an ancestry to a host hole, possibly
the replaced hole. An inserted starting edge has an ancestry to an inserted
hole. The only problematic crossing is from insertion back to host: it
would create a strict ancestry from the replaced hole to a host hole.
Binding-freedom excludes this because the replaced hole has a source.
An inserted-to-inserted adjacency reflects to the original insertion using
the previously proved absence of a self-loop at the replaced hole.

For shallow plugging closure, split the first edge and the first successor
of an alleged forbidden path by host/inserted origin. Host-to-host and
inserted-to-inserted paths contradict the respective original restriction.
A host-to-inserted crossing produces a forbidden host path to the replaced
hole. An inserted-to-host crossing contradicts the host restriction at that
sourced hole. Source witnesses survive the incidence-list maps.

\lean{Copying/BindingClosure}{14}{21}{BindingClosure.adjacent_host_insert}


\lean{Copying/BindingClosure}{22}{32}{BindingClosure.adjacent_insert_host}


\lean{Copying/BindingClosure}{33}{68}{BindingClosure.ancestor_origin}


\lean{Copying/BindingClosure}{69}{105}{BindingClosure.shallow_plug}


The addressed sourced-hole predicate descends through boxes to the
selected shallow occurrence. Shallow closure of the recursive path predicate
uses the path theorem at the current depth and retains the body predicates
of all surviving or inserted boxes. Replacing only a box body changes no
outer incidence and no outer binding/hole classification; the pullback
lemma therefore preserves the outer restriction.

Induction on the addressed occurrence combines these two cases and proves
closure at arbitrary nesting depth. The final binding-free plugging theorem
adds the existing global focus-freedom closure. No copying transition is
assumed or encoded by this argument.

\lean{Copying/BindingClosure}{106}{110}{BindingClosure.SourcedHole}


\lean{Copying/BindingClosure}{111}{125}{BindingClosure.paths_shallow}


\lean{Copying/BindingClosure}{126}{144}{BindingClosure.shallow_replaceBox}


\lean{Copying/BindingClosure}{145}{168}{BindingClosure.paths_plugAt}


\lean{Copying/BindingClosure}{169}{175}{BindingClosure.plugAt}


Recursive label equivalence preserves both binding-label and hole-label
classification: primitive labels agree, boxes match boxes, and mixed cases
are impossible. A recursive graph isomorphism preserves adjacency and maps
source witnesses by its vertex bijection. The pullback lemma gives shallow
invariance. Induction on depth then applies the same argument to corresponding
box bodies. Applying it to an isomorphism and its inverse, together with
focus-freedom invariance, proves preservation and reflection of binding-freedom.

\lean{Copying/BindingInvariance}{8}{20}{BindingInvariance.bindingLabel_equiv}


\lean{Copying/BindingInvariance}{21}{33}{BindingInvariance.holeLabel_equiv}


\lean{Copying/BindingInvariance}{34}{42}{BindingInvariance.shallow_iso}


\lean{Copying/BindingInvariance}{43}{64}{BindingInvariance.paths_iso}


\lean{Copying/BindingInvariance}{65}{72}{BindingInvariance.iso_iff}


\subsection{Term-hole checks and a zero-input counterexample}
\termdef{copying-zero-input-counterexample}


Every already constructed term translation is binding-free because it is
focus-free and hole-free. This is a theorem about terms, not yet about
lambda-context translations. A stand-alone hole corolla is binding-free:
its targets are outputs and its sources are inputs, so no positive return
adjacency is possible. Finally, the hole assigned to any term judgement has
interface $\star\Rightarrow\star^k\otimes\diamond^h$ and hence has the
source witness required by the plugging theorem, for every $k$ and $h$.

\lean{Copying/BindingExamples}{11}{14}{BindingExamples.translated_bindingFree}


\lean{Copying/BindingExamples}{15}{29}{BindingExamples.hole_bindingFree}


\lean{Copying/BindingExamples}{30}{35}{BindingExamples.term_hole_sourced}


Consider a context with two holes in series,
\[
 \chi_0: \epsilon\Rightarrow\diamond,
 \qquad \chi_1:\diamond\Rightarrow\epsilon.
\]
Their names are distinct. The context is binding-free under the literal
source-based definition: $\chi_0$ has no source, while $\chi_1$ has no target.
Now replace $\chi_0$ with a term weakening followed by an instance:
\[
 \epsilon\xrightarrow{\mathrm{weakening}_\star}\star
 \xrightarrow{\mathsf I}\diamond.
\]
This insertion is focus-free and hole-free, hence binding-free. After plugging,
there is a one-edge path from the source of $\mathsf I$ to the source of
$\chi_1$. The result is not binding-free.

Lean checks both finite constructions, both original binding-freedom proofs,
the uniqueness of the host's hole labels, the resulting forbidden path, and
the negation of unrestricted hole-replacement closure. The source witness
and the final adjacency are explicit vertices of the constructed graph;
no external decision procedure or new axiom is used.

This exposes a boundary in the literal general formulation of Lemma~E.1.
Its proof prefixes the old hole to a path, but that prefix has no starting
source if the hole has no inputs. The qualified theorem proved here suffices
for plugging at translated term holes. It does not silently strengthen
Definition~11.2, nor remove zero-input holes from the observation class.
The broader correspondence audit must record this qualification.

\lean{Copying/BindingExamples}{36}{38}{BindingExamples.zeroInputContext}


\lean{Copying/BindingExamples}{39}{41}{BindingExamples.zeroInputReplacement}


\lean{Copying/BindingExamples}{42}{53}{BindingExamples.zeroInputContext_bindingFree}


\lean{Copying/BindingExamples}{54}{60}{BindingExamples.zeroInputReplacement_bindingFree}


\lean{Copying/BindingExamples}{61}{68}{BindingExamples.zeroInputContext_unique_holes}


\lean{Copying/BindingExamples}{69}{82}{BindingExamples.zeroInput_plug_not_bindingFree}


\lean{Copying/BindingExamples}{83}{93}{BindingExamples.unrestricted_hole_closure_false}


The full build and all three axiom audits pass. Remaining work in
Appendix~B.1, item~3 includes constructing copying observations, translating
lambda-contexts and proving their binding-freedom and plugging correspondence,
and deriving term-context congruence. That item is not marked complete.
